\documentclass[10pt,a4paper]{amsart}
\usepackage[margin=19mm]{geometry}
\usepackage{amsmath,amssymb,mathtools}
\usepackage[scr=rsfs,cal=euler]{mathalfa}
\usepackage{xfrac}
\usepackage[unicode,hidelinks,bookmarks=false]{hyperref}
\newcommand{\ie}{i.e.\ }
\newcommand{\cf}{cf.\ }
\newcommand{\resp}{respectively\ }
\newcommand{\Subsection}{\subsection}
\providecommand{\nobreakdash}{\nobreak\hbox{-}}
\setlength{\emergencystretch}{2em}
\multlinegap0pt
\usepackage{bm}
\usepackage{bbm}
\usepackage{dsfont}
\usepackage{xspace}
\usepackage{cancel}
\usepackage{mathtools}
\usepackage{amsthm}
\makeatletter
\@ifundefined{satz}{\newtheorem{satz}{Satz}}{}
\@ifundefined{lemma}{\newtheorem{lemma}[satz]{Lemma}}{}
\@ifundefined{theorem}{\newtheorem{theorem}{Theorem}}{}
\makeatother
\newcommand{\RR}{\mathbb R}
\newcommand{\norm}[1]{\left\| #1 \right\|}
\newcommand{\pa}{\partial}
\title[Derivation and well-posedness for asymptotic models of cold ]{Derivation and well-posedness for asymptotic models of cold plasmas}
\author{Diego Alonso-Or\'an, \'Angel Dur\'an, Rafael Granero-Belinch\'on}
\date{Source: arXiv:2305.13922v2 (CC BY 4.0)}
\begin{document}
\maketitle

\noindent\emph{Source and licence.} Derivation and well-posedness for asymptotic models of cold plasmas. Version arXiv:2305.13922v2 (CC BY 4.0), distributed under the Creative Commons Attribution 4.0 International licence, as recorded in the arXiv licence metadata for this exact version and shown on the abstract page https://arxiv.org/abs/2305.13922v2. Full rights notice: licenses/arxiv-2305.13922-CC-BY-4.0.txt.\par\smallskip

\noindent\textit{Editorial scope.} Counted unit: the Theorem whose source label is \texttt{theoremBouss} in the article (source0.tex lines 691--695, source label \texttt{theoremBouss}), delivered together with the article's own complete original proof of it (source0.tex lines 697--823, one \texttt{proof} environment of 127 lines). No mathematics is written, repaired or re-derived: statement and proof are the article's own text line for line. Required context retained uncounted: eq:Boussinesq (lines 190--195); adg4 (lines 202--205); comp:helmholtz (lines 289--291); lema:LtoL2 (lines 659--664); final:eq:uni:2 (lines 1012--1014); theorem1 (lines 1016--1018). Every definition, display and label of those lines is retained unchanged. Omitted from this package and not redistributed: 1--189, 196--201, 206--288, 292--658, 665--690, 696--696, 824--1011, 1015--1015, 1019--1609. Nothing inside a retained excerpt is omitted or altered: every retained body is delivered exactly as the article's lines are written, so no omission receipt of any kind accompanies this package. Citations: the retained excerpts carry no citation command at all, so no bibliography entry of the article is transcribed and no reference is redistributed. Reference adaptation: none was needed and none was made; every reference inside the delivered unit resolves inside it, so no unresolved reference or citation is produced. Printed slips kept literal: no printed word, symbol, hypothesis or number of the counted statement or of its proof was altered. Layout and class: the article's own class is \texttt{amsart} with packages including stmaryrd, showlabels, babel, pst-grad, pst-plot, pstricks, amsmath, amssymb, mathrsfs, amsthm, mathtools, tikz, mathcomp, wasysym; the wrapper is the lane's pinned \texttt{amsart} editorial wrapper at 10pt on A4, which restates the theorem-like environments the retained text needs and adds the article's own macro definitions that the retained text needs (\texttt{\textbackslash RR}, \texttt{\textbackslash norm}, \texttt{\textbackslash pa}), with extra wrapper packages (bm, bbm, dsfont, xspace, cancel, mathtools). The delivered numbering is the unit's own. Assets: no retained span contains a figure environment, a graphics-inclusion command, a PDF-page inclusion, a table or a TikZ picture; no image file is copied into this package, the package declares no asset dependency and no image file is redistributed with it. Licence: version arXiv:2305.13922v2 of this article is distributed under the Creative Commons Attribution 4.0 International licence, as recorded in the arXiv licence metadata for this exact version and shown on the abstract page https://arxiv.org/abs/2305.13922v2. A full rights notice is delivered at licenses/arxiv-2305.13922-CC-BY-4.0.txt. Characters: every retained excerpt is pure ASCII, so the delivered engine, XeLaTeX with its default Latin Modern text font, reports no missing character.
\begin{subequations}\label{eq:Boussinesq}
\begin{align}
h_t+(hv)_x+v_x&=0,\\
v_t+vv_x+\left[\mathscr{L},\mathscr{N}h\right]h+\mathscr{N}h&=0,
\end{align}
\end{subequations}

\begin{equation}\label{adg4}
\widehat{\mathscr{L}h}(\xi)=\frac{\xi^2}{1+\xi^2}\hat{h}(\xi),
\; \widehat{\mathscr{N}h}(\xi)=\frac{i\xi}{1+\xi^2}\hat{h}(\xi).
\end{equation}

\begin{equation}\label{comp:helmholtz}
\pa_{x}^{2}\mathscr{Q}f(x)=\left( \mathscr{Q}-\text{I}\right)f(x),\; x\in\mathbb{R},
\end{equation}

\begin{lemma}\label{lema:LtoL2}
Let  $k\in \mathbb{N}$ and $f\in L^{2}(\RR)$ be smooth enough. Then, there exists a constant $C>0$ such that 
\begin{equation}\label{estimate:lema:LtoL2}
\norm{\pa_{x}^{k}f}_{L^{2}} \leq C\left(\norm{f}_{L^{2}}+\norm{\sqrt{\mathscr{L}}\pa_{x}^{k}f}_{L^{2}}\right).
\end{equation}
\end{lemma}

\begin{equation}\label{final:eq:uni:2}
h_{t}=-\frac{1}{2}\left(3 hh_{x}+[\mathscr{Q},\mathscr{Q}h_{x}]h - \mathscr{Q}h_{x}-h_{x} \right). 
\end{equation}

\begin{theorem}\label{theorem1}
Let $s>\frac{3}{2}$  and $h_0\in H^{s}(\RR)$ {with mean zero}. Then there exist $T_{\max}>0$ and a unique local solution $h\in C([0,T_{\max}),H^{s}(\RR))$ of \eqref{final:eq:uni:2} with $h(0)=h_{0}$.
\end{theorem}

\begin{theorem}\label{theoremBouss}
For $(h_0,v_0)\in H^2\times H^3$  there exist a time $0<T_{max}$ and a unique solution 
\[ (h,v)\in C((0,T_{max}), H^2\times H^3)\]
of the ivp of \eqref{eq:Boussinesq} with $h(x,0)=h_{0}(x), v(x,0)=v_{0}(x), x\in\mathbb{R}$.
\end{theorem}

\begin{proof}
We first focus on obtaining some a priori estimates.
%, since the construction of the solutions can be done by standard approximation arguments. Instead of using the standard norm in $H^2\times H^3$, we are going to use the norm in $\mathcal{X}$.\textcolor{red}{Son el mismo espacio pero diferente norma. No se si escribir esto asi o olvidar la notacion $\mathcal{X}$} We will roughly sketch the approximation approach at the end of the proof. 
We define the energy
\begin{equation}\label{energy:def}
\mathcal{E}(t)=\norm{(h,v)}_{L^{2}\times L^2}^{2} 	+\norm{\sqrt{\mathscr{L}}\pa_{x}^{2}h}_{L^{2}}^{2}+ \norm{\pa_{x}^{3}v}_{L^{2}}^{2},
\end{equation}
where the norm considered in the space $L^{2}(\mathbb{R})\times L^{2}(\mathbb{R})$ is given by
$$\norm{(f,g)}_{L^{2}\times L^2}=\left(||f||_{L^{2}}^{2}+||g||_{L^{2}}^{2}\right)^{1/2}, f,g\in L^{2}(\mathbb{R}).$$
Multiplying the first equation of \eqref{eq:Boussinesq} by $h$, the second by $v$, adding the resulting equalities and integrating on $\mathbb{R}$ we have
\begin{align*}
\frac{1}{2}\frac{d}{dt}\norm{(h,v)}_{L^{2}\times L^2}^{2}=-\int_{\RR} \left((hv)_{x}+v_{x}\right) h \  dx- \int_{\RR} \left(vv_{x}+[\mathscr{L}, \mathscr{N}h]h+\mathscr{N}h \right) v \  dx.
\end{align*}
Now integration by parts, the application of H\"older and Young  inequalities, and the Sobolev embedding $H^{\frac{1}{2}+\epsilon}(\RR)\hookrightarrow L^{\infty}(\RR)$ for $\epsilon>0$ lead to 
\begin{align}\label{L2:estimate:system3}
\frac{1}{2}\frac{d}{dt}\norm{(h,v)}_{L^{2}\times L^2}^{2}\leq C \norm{v_x}_{L^{\infty}}\norm{h}^{2}_{L^{2}} \leq \norm{\pa_{x}^{3}v}^{3}_{L^2}+\norm{h}_{L^{2}}^{3}.
\end{align}
Next, we deal with the other terms in \eqref{energy:def}. Multiplying the first equation in \eqref{eq:Boussinesq} by $ \mathscr{L}\partial_{x}^{4}h$ and integrating we have 
\begin{align}\label{eq:Lh}
\int_{\RR} \mathscr{L}\pa_{x}^{4}h h_{t} \  dx=-\int_{\RR}(hv)_{x} \mathscr{L}\pa_{x}^{4}h \ dx - \int_{\RR} v_{x} \mathscr{L}\pa_{x}^{4}h \  dx.
\end{align}
On the other hand, multiplying the second equation in  \eqref{eq:Boussinesq} by $\pa_{x}^{6}v$ and integrating we obtain
\begin{align}\label{eq:H3v}
\int_{\RR} v_{t} \pa_{x}^{6}v\  dx &=-\int_{\RR}vv_{x} \pa_{x}^{6}v \ dx - \int_{\RR}[\mathscr{L}, \mathscr{N}h] h \pa_{x}^{6}v  \ dx - \int_{\RR} \mathscr{N}h \pa_{x}^{6} v \  dx .
\end{align}
Since $-\mathscr{L}=\pa_{x}\mathscr{N}$, integrating by parts we find that the last term in \eqref{eq:Lh} is given by
\begin{equation}\label{rewrittingI2}
 - \int_{\RR} v_{x} \mathscr{L}\pa_{x}^{4}h \  dx=\int_{\RR} v_{x} \mathscr{N}\pa_{x}^{5}h \  dx= -\int_{\RR} \pa_{x}^{6} v \mathscr{N}h \ dx .   
 \end{equation}
Moreover, we have that
\begin{equation}\label{rewritingtemporal}
\int_{\RR} h_{t}\mathscr{L}\pa_{x}^{4}h  \  dx=\frac{1}{2}\frac{d}{dt}\norm{\sqrt{\mathscr{L}}\pa_{x}^{2}h} ^{2}_{L^{2}}, \quad  -\int_{\RR} v_{t} \pa_{x}^{6}v\  dx=  \frac{1}{2}\frac{d}{dt}\norm{\pa_{x}^{3}v}^{2}_{L^{2}}.
\end{equation}
Then, adding \eqref{eq:Lh} and \eqref{eq:H3v}, and using \eqref{rewrittingI2},\eqref{rewritingtemporal} lead to
\begin{equation}\label{adg13}
\frac{1}{2}\frac{d}{dt}\left( \norm{\sqrt{\mathscr{L}}\pa_{x}^{2}h}_{L^{2}}^{2}+ \norm{\pa_{x}^{3}v}_{L^{2}}^{2} \right) =\underbrace{ -\int_{\RR}(hv)_{x} \mathscr{L}\pa_{x}^{4}h \ dx}_{I_{1}}+\underbrace{\int_{\RR}vv_{x} \pa_{x}^{6}v \ dx}_{I_{2}}+\underbrace{\int_{\RR}[\mathscr{L}, \mathscr{N}h] h \pa_{x}^{6}v  \ dx}_{I_{3}}.
\end{equation}
We now estimate each of the integrals $I_{i}$ in (\ref{adg13}).
First, notice that integration by parts yields 
\begin{equation*}
I_{2}=-\int_{\RR} \pa_{x}^{2}(vv_{x})\pa_{x}^{3}v  \  dx =-\frac{3}{2}\int_{\RR}(\pa_{x}^{3}v)^{2} v_{x} \  dx,
\end{equation*}
and thus
\begin{align}\label{est:I2}
|I_{2}|\leq C  \norm{\pa_{x}v}_{L^{\infty}}\norm{\pa_{x}^{3}v}^{2}_{L^{2}}\leq  C \norm{\pa_{x}^{3}v}^{3}_{L^{2}},
\end{align}
where in the second inequality we used the Sobolev embedding $H^{\frac{1}{2}+\epsilon}(\RR)\hookrightarrow L^{\infty}(\RR)$ for $\epsilon>0.$

Integrating by parts the first term $I_{1}$ we find that
\begin{equation*}
I_{1}=-\int_{\RR}\pa_{x}^{3}(hv)\pa_{x}^{2}\mathscr{L}h \ dx =\underbrace{-\int_{\RR}\left( \pa_{x}^{3}v h+ \pa_{x}^{2}vh_{x}+v_{x}\pa_{x}^{2}h \right)\pa_{x}^{2}\mathscr{L}h \ dx}_{J_{1}}\underbrace{-\int_{\RR}v\pa_{x}^{3}h \pa_{x}^{2}\mathscr{L}h \ dx}_{J_{2}} .
\end{equation*}
We use H\"older inequality to estimate $J_{1}$ as
\begin{align}\label{est:J1}
|J_{1}|&\leq C \left( \norm{h}_{L^{\infty}}\norm{\pa_{x}^{3}v}_{L^{2}}+\norm{\pa_{x}^{2}v}_{L^{\infty}}\norm{\pa_{x}h}_{L^{2}}+\norm{v_{x}}_{L^{\infty}}\norm{h_{xx}}_{L^{2}}\right)\norm{\pa_{x}^{2}\mathscr{L}h}_{L^{2}} \nonumber \\
& \leq C \norm{h_{xx}}_{L^{2}}\norm{\pa_{x}^{3}v}_{L^{2}}\norm{\pa_{x}^{2}\mathscr{L}h}_{L^{2}}.
\end{align}
As for $J_{2}$, using the identity (cf. \eqref{comp:helmholtz})
\begin{equation}\label{adg14}
\mathscr{L}=-\partial_{x}^2 (1-\pa_{x}^{2})^{-1}= \mathrm{Id}-\mathscr{Q},
\end{equation}
we obtain that
\begin{equation*}\label{rewrite:J2}
J_{2}=-\int_{\RR}v \pa_{x}^{3}h \left(\pa_{x}^{2}h-\mathscr{Q}\pa_{x}^{2}h\right) \ dx =- \frac{1}{2}\int_{\RR}v \pa_{x}(\pa_{x}^{2}h)^{2} \ dx + \int_{\RR}v \pa_{x}^{3}h \mathscr{Q}\pa_{x}^{2}h \ dx.
\end{equation*}
Using integration by parts in both terms we infer that
\begin{equation}\label{est:J2}
|J_{2}|\leq C \norm{v_{x}}_{L^{\infty}} \norm{h_{xx}}^{2}_{L^{2}}\leq \norm{\pa_{x}^{3}v}_{L^{2}} \norm{h_{xx}}^{2}_{L^{2}}.
\end{equation}
In a similar way, expanding the commutator and using (\ref{adg14}) again lead to
\begin{align*}
I_{3}=\int_{\RR} \left( \mathscr{L}(\mathscr{N}h h)-\mathscr{N}h\mathscr{L}h\right)\pa_{x}^{6}v \  dx=\underbrace{- \int_{\RR} \mathscr{Q}(\mathscr{N}h h)\pa_{x}^{6}v \  dx}_{K_{1}} +\underbrace{ \int_{\RR} \mathscr{N}h  \mathscr{Q}h \pa_{x}^{6}v \  dx}_{K_{2}}.
\end{align*}
We rewrite both terms $K_{1}$ and $K_{2}$ as
\begin{align*}
K_{1}= -\int_{\RR} (\mathscr{N}hh)_{x} \mathscr{L}\pa_{x}^{3}v \  dx,
\end{align*}
and
\begin{align*}
K_{2}=-\int_{\RR} (\mathscr{N}h\mathscr{Q}h)_{xxx} \pa_{x}^{3}v \  dx &=- \frac{1}{2} \int_{\RR} \pa_{x}^{4}\left( (\mathscr{Q}h)^{2}
\right) \pa_{x}^{3}v \  dx =-\int_{\RR}\pa_{x}^{2}\left( (\mathscr{Q}h_{x})^{2}-\mathscr{Q}h\mathscr{L}h\right) \pa_{x}^{3}v \  dx.
\end{align*}
Using the fact that, cf. (\ref{adg4}), 
\begin{equation}\label{adg15}
\norm{\sqrt{\mathscr{L}}f}_{L^2}\leq C \norm{f}_{L^{2}},\; f\in L^{2},
\end{equation}
then the term $K_{1}$ can be bounded by
\begin{equation}\label{bound:K1prev}
|K_{1}|\leq C \norm{h_{xx}}_{L^{2}}^{2}\norm{\mathscr{L}\pa_{x}^{3}v}_{L^{2}} \leq \norm{h_{xx}}_{L^{2}}^{2}\norm{\pa_{x}^{3}v}_{L^{2}},
\end{equation}
for some constant $C$. As far as $K_{2}$ is concerned, by expanding the derivatives, tedious but a straightforward computation and using (\ref{adg15}) again shows that
%\begin{equation}\label{rewrite:K2}
%K_{2}=\int_{\RR} \mathscr{Q}h\mathscr{L}\pa_{x}^{2}h\pa_{x}^{3}v \  dx + \texttt{l.o.t}= \int_{\RR} \sqrt{\mathscr{L}}\left(\mathscr{Q}h   \pa_{x}^{3}v\right) \sqrt{\mathscr{L}}\pa_{x}^{2}h \  dx + \texttt{l.o.t},
%\end{equation}
%and hence
\begin{equation}\label{bound:K2prev}
|K_2| \leq C \norm{\mathscr{Q}h}_{L^{\infty}} \norm{\pa_{x}^{3}v}_{L^{2}}\norm{\sqrt{\mathscr{L}}\pa_{x}^2 h}_{L^{2}} \leq C    \norm{\pa_{x}^{3}v}_{L^{2}}\norm{h_{xx}}_{L^{2}} \norm{\sqrt{\mathscr{L}}\pa_{x}^2 h}_{L^{2}},
\end{equation}
for some constant $C$. From \eqref{est:I2}, \eqref{est:J1}, \eqref{est:J2}, \eqref{bound:K1prev}, and \eqref{bound:K2prev} we conclude that

%Therefore, collecting the previous estimates \eqref{est:I2}, \eqref{est:J1}, \eqref{est:J2}, \eqref{bound:K1prev} and \eqref{bound:K2prev} we conclude that
\begin{align}
\frac{1}{2}\frac{d}{dt}\left( \norm{\sqrt{\mathscr{L}}\pa_{x}^{2}h}_{L^{2}}^{2}+ \norm{\pa_{x}^{3}v}_{L^{2}}^{2} \right) &\leq  C \bigg( \norm{\pa_{x}^{3}v}^{3}_{L^{2}}+\norm{\pa_{x}^{3}v}_{L^{2}} \norm{h_{xx}}^{2}_{L^{2}} \nonumber \\
&\hspace{3cm} +\norm{\pa_{x}^{3}v}_{L^{2}}\norm{h_{xx}}_{L^{2}} \norm{\sqrt{\mathscr{L}}\pa_{x}^2 h}_{L^{2}} \bigg).\label{adg16}
\end{align}

Applying Lemma \ref{lema:LtoL2} and Young's inequality to (\ref{adg16}) leads to
\begin{align}\label{H2:estimate:system3}
\frac{1}{2}\frac{d}{dt}\left( \norm{\sqrt{\mathscr{L}}\pa_{x}^{2}h}_{L^{2}}^{2}+ \norm{\pa_{x}^{3}v}_{L^{2}}^{2} \right) \leq  C \left( \norm{\sqrt{\mathscr{L}}\pa_{x}^2 h}^{3}_{L^{2}}+\norm{h}^{3}_{L^{2}} +\norm{\pa_{x}^{3}v}^{3}_{L^{2}}\right),
\end{align}
for some constant $C$. 
Combining estimates \eqref{L2:estimate:system3} and \eqref{H2:estimate:system3} yields
\begin{equation}\label{energy:estimate:bouss:nolocal}
\mathcal{E}'(t)\leq C \mathcal{E}^{\frac{3}{2}}(t),
\end{equation}
which ensures a local time of existence $t^{*}>0$ such that $\mathcal{E}(t)\leq 4 \mathcal{E}(0), 0<t<t^{*}.$
In order to construct the solutions, we first define the approximate problems using mollifiers (cf. proof of Theorem \ref{theorem1}). More precisely, the regularized system is given by
\begin{subequations}
	\begin{align}
		h^{\epsilon}_t+\mathcal{J}_{\epsilon}(\mathcal{J}_{\epsilon}h^{\epsilon}\mathcal{J}_{\epsilon}v^{\epsilon})_x+\mathcal{J}_{\epsilon}\mathcal{J}_{\epsilon}v^{\epsilon}_x&=0, \label{regularized:eq:Boussinesq} \\
		v^{\epsilon}_t+\mathcal{J}_{\epsilon}(\mathcal{J}_{\epsilon}v^{\epsilon}\partial_{x}\mathcal{J}_{\epsilon}v^{\epsilon})+\left[\mathscr{L},\mathscr{N}h^{\epsilon}\right]h^{\epsilon}+\mathscr{N}h^{\epsilon}&=0.\label{regularized:eq:Boussinesq2}
	\end{align}
\end{subequations}


 By the properties of the mollifiers we can repeat the previous energy estimates and provide the same a priori bounds for the regularized system of \eqref{regularized:eq:Boussinesq}-\eqref{regularized:eq:Boussinesq2}. Hence, we will find a uniform time of existence $T_{max}>0$ for the sequence of regularized problems. To conclude the argument, we pass to the limit. Furthermore, the continuity in time for the solution (instead of merely weak continuity) is obtained as follows: first, the energy estimate \eqref{energy:estimate:bouss:nolocal} yields the strong right continuity at $t=0$. Moreover, it is easy to check that changing variables $\tilde{t}=-t$ provides the strong left continuity at $t=0$ and hence the continuity in time for the solutions. To conclude let us remark that the uniqueness follows by a classical contradiction argument as in Theorem \ref{theorem1}.
\end{proof}

\end{document}
